% Lösungen Einheit 4 von 8 – Ziehen ohne Zurücklegen und Umlegen (zusammenbau v0.1)
\einheitenkopf{Einheit 4 von 8 · Ziehen ohne Zurücklegen und Umlegen}

\erg{21}{a) p ändert sich \quad b) $\frac{10}{25} \cdot \frac{9}{24} = \frac{3}{20}$ \quad c) P(A) $= \frac{16}{40} \cdot \frac{15}{39} \cdot \frac{14}{38} \approx 0{,}0567$; P(B) $= 1 - \frac{24}{40} \cdot \frac{23}{39} \cdot \frac{22}{38} \approx 0{,}795$ \quad d) $\frac{8}{32} \cdot \frac{7}{31} \cdot \frac{6}{30} \cdot \frac{5}{29} \approx 0{,}00195$ \quad e) $\frac{3}{4} \cdot \frac{3}{5} + \frac{1}{4} \cdot \frac{2}{5} = \frac{11}{20}$ \quad f) Gleich groß, je $\frac{3}{7}$: Die vier Pfade zum dritten Zug ergeben zusammen $\frac{3}{7}$ – oder Symmetrie: jede Ziehungsposition ist gleich gut.}

21c)

\baumdrei{V/\frac{16}{40},N/\frac{24}{40}}{V/\frac{15}{39},N/\frac{24}{39}}{V/\frac{16}{39},N/\frac{23}{39}}{V/\frac{14}{38},N/\frac{24}{38}}{V/\frac{15}{38},N/\frac{23}{38}}{V/\frac{15}{38},N/\frac{23}{38}}{V/\frac{16}{38},N/\frac{22}{38}}

\erg{22}{a) Nils übersieht die umgelegte Kugel: B hat danach 5 Kugeln, je nach Farbe 2 oder 1 rote: $\frac{1}{2} \cdot \frac{2}{5} + \frac{1}{2} \cdot \frac{1}{5} = \frac{3}{10}$ \quad b) Sie gilt nur unter der Bedingung, was im ersten Zug gezogen wurde – die Urne hat dann eine Kugel weniger und je nach Farbe eine andere Zusammensetzung. \quad c) $1 - \frac{22}{24} \cdot \frac{21}{23} \cdot \frac{20}{22} \approx 0{,}239$ \quad d) $\frac{3}{7} \cdot \frac{4}{6} + \frac{4}{7} \cdot \frac{3}{6} = \frac{4}{7}$}
